\documentclass[10pt]{article}
\usepackage[margin=20mm]{geometry}
\usepackage{amsmath,amssymb,fontspec,unicode-math,xurl,hyperref}
\IfFontExistsTF{DejaVu Serif}{\setmainfont{DejaVu Serif}}{\setmainfont{Latin Modern Roman}}
\setmathfont{latinmodern-math.otf}
\setlength{\parindent}{0pt}\setlength{\parskip}{6pt}\sloppy
\begin{document}
\section*{Scientific verification report: I25-40117}
\section*{1. Identity and exact reviewed version}
Paper: Position: Principles of Animal Cognition to Improve LLM Evaluations. Authors: Sunayana Rane, Cyrus F. Kirkman, Graham Todd, Amanda Royka, Ryan M. C. Law, Erica A. Cartmill, and Jacob G. Foster. ICML 2025, PMLR 267:82051-82061. Assigned display number: 0018.\par
The reviewed source is the 11-page final PMLR PDF, SHA-256 461ff952375293337551d22b9528546048b186ddb829f43a0e3bde9bdc0cb34b, downloaded at 2026-10-24T04:31:04.670160+00:00. Page references are 1-based PDF pages and match printed internal pages. No appendix is present. The public preliminary drafts are not substituted for this version. [Rane-PMLR (2025)]\par
Review date: 24 October 2026 (UTC).\par
\section*{2. Summary and supported judgment}
Supported outcome: Minor. One concrete reporting incompatibility is demonstrated: the Experiment 2 endpoint-exclusion rule permits test distances 2-4, whereas the results give 2-6. Which description matches the implemented tests remains unresolved. The numerical chance-performance and uncertainty claims require additional evidence for independent verification; no error in the reported accuracy of 0.40 or SD of 0.50 has been established.\par
The five proposed evaluation principles remain defensible. The case study supports examining sensitivity to task formulation. It does not, by itself, isolate a unique reasoning mechanism. The paper’s own distinction between behavior and mechanism and its contextual discussion are important counterevidence against reading it as a blanket denial of all transitive competence. No Major issue is supported.\par
\section*{3. Definitions, coverage, and evidence limits}
The empirical question concerns transitive inference (TI): extending an ordinal relation to pairs not explicitly compared. Experiment 1 states ordinal premises in a prompt and varies operators and stimuli. Experiment 2 supplies correct/incorrect feedback on adjacent binary choices without explicitly stating a total-order assumption. These tasks involve overlapping but different demands. Symbolic distance is rank separation: adjacent elements have distance 1. [Rane-PMLR (2025), pp.5-8]\par
All 11 source pages were read, and Figures 1-2 and Table 1 were visually checked against the exact original. The check included quantifiers, training/test denominators, endpoint exclusions, figure signs, illustrative alternative mechanisms, source identity, and public version/materials searches. It did not reproduce GPT-4o outputs or validate every cited animal experiment.\par
The PMLR related-material section links the PDF and OpenReview, with no separate software or supplement link. No experiment repository or raw-response package was located in the inspected PDF, author publication links, or targeted searches. This is a bounded availability result, not proof that such records do not exist. OpenReview returned a browser-verification page, leaving its discussion and attachments unverified.\par
\section*{4. Problem 1: the stated test rule and distance range cannot both hold}
\subsection*{Original claim and suspect step}
On p.7, section 6.2 describes seven ordered stimuli A through G and novel test pairs that exclude the terminal elements. Section 7.3 on the same page reports tests at symbolic distances 2-6. The distance definition in Figure 2, p.8, is rank separation. [Rane-PMLR (2025)]\par
\subsection*{Complete finite check}
After removing A and G, the available ranks are B, C, D, E, F. Enumerating all unordered nonadjacent pairs gives exactly:\par
Pair | Distance\par
B,D | 2\par
B,E | 3\par
B,F | 4\par
C,E | 2\par
C,F | 3\par
D,F | 2\par
Thus the admissible distance set is \{2,3,4\}, with maximum 4. Distance 5 requires A,F or B,G; distance 6 requires A,G. All of those pairs contain a terminal element. Reversing the display order of a pair does not change its rank separation. The finite enumeration was independently executed using the standard library and saved in verification/0018/independent\_results.json.\par
\subsection*{Counterevidence and alternative readings}
The reported range could be a carryover from Experiment 1, which includes all pairs. Alternatively, the actual Experiment 2 tests may have included endpoints. The abbreviated Figure 1 schematic is not a trial log and cannot decide between these possibilities.\par
A nearby statement assigning B an always-incorrect status is also inconsistent with B being an interior member of the adjacent chain; it plausibly intends G. This is treated as related descriptive ambiguity, not as a second independent defect or evidence that the experiment used incorrect labels.\par
\subsection*{Impact and conditional repair}
Judgment: Minor, a bounded methods/results incompatibility. Endpoint exclusion matters because endpoints have uniformly reinforced status, so choices involving them can be answered without general nonadjacent inference. The inconsistency prevents a precise reconstruction of the intended control. It does not show that the reported 0.40 accuracy is false.\par
If endpoint exclusion describes the implementation, the admissible range is 2-4; this conditional mathematical correction is verified. If distances 5 or 6 were actually tested, the method description and endpoint-free interpretation require revision. The actual pair list must determine the repair, followed by recalculation of any affected summaries. No empirical repair has been verified or adopted.\par
\section*{5. Problem 2: chance-performance and uncertainty claims cannot yet be independently verified}
\subsection*{Claim, missing information, and precise uncertainty}
Section 7.3 characterizes Experiment 2 performance as at chance and reports no distance differences. Table 1 gives accuracy 0.40 and SD 0.50. Figure 2 presents 95\% intervals and fitted lines for Experiment 1. The text specifies ten task iterations and sixty training trials per iteration, but does not establish the independent test-response denominator or the SD’s aggregation unit. The inspected materials also do not identify a reproducible model snapshot, inference settings, complete response history, or sufficient GLM family/link and uncertainty details. [Rane-PMLR (2025), pp.6-8]\par
Ten iterations cannot automatically be equated with ten test observations, and the six hundred training presentations are not test responses. An SD near 0.50 can be compatible with binary outcomes near accuracy 0.4; it is not evidence of an arithmetic error. The Experiment 2 paragraph’s pointer to Figure 2 does not supply its missing trial data because that figure describes Experiment 1.\par
\subsection*{Independent check of what accuracy alone can establish}
To illustrate dependence on the missing denominator, consider hypothetical independent Bernoulli responses under a known null probability 0.5. An exact two-sided binomial test sums probabilities no greater than the observed probability. Two hypothetical samples have identical accuracy 0.40:\par
Hypothetical sample | Accuracy | Exact two-sided p-value\par
4 correct of 10 | 0.40 | 0.75390625\par
400 correct of 1000 | 0.40 | 2.7284641560660184e-10\par
These are independently calculated illustrations, not reconstructions of the paper’s data or its SD. They prove that an accuracy value alone does not determine its statistical relation to chance. Actual observations may also be clustered within conversations or task instances, so an independent-binomial calculation could be inappropriate even after obtaining a raw count.\par
Failure to reject a chance null is not, by itself, an equivalence result. A formal equivalence claim would additionally require a justified tolerance and suitable uncertainty calculation.\par
\subsection*{Counterevidence, impact, and repair status}
Judgment: Needs Evidence for independent statistical verification. The authors may intend a descriptive characterization of 40\% as roughly near a 50\% baseline; that reading is acknowledged and is not disproved. Likewise, all six displayed fitted slopes are visually positive, consistent with their sign claim. No contrary numerical result is asserted.\par
The unresolved material is needed to assess sampling uncertainty, distance effects, and reproducibility. Release the actual test pairs and responses, independent units, aggregation rules, model/inference metadata, and uncertainty procedure. Reanalysis should respect dependence and distinguish estimation, a difference test, and equivalence. The published estimates, intervals, and fitted coefficients have not been independently reproduced.\par
\section*{6. Mechanism and scope checks: no additional demonstrated defect}
\subsection*{Inductive transfer versus deduction from declared ordinal premises}
Figure 1C asks the system to learn correct word choices through feedback; it does not announce that the hidden relation is transitive. A concrete countermodel shows the resulting inference boundary. Let one rule choose the earlier item in A>B>C>D>E>F>G. Let a second rule agree everywhere except for the unseen B,D pair, where it chooses D. Both rules make identical predictions for all six trained adjacent pairs, and hence for all sixty repetitions, but they disagree on B,D. The second rule has a cycle B>C>D>B.\par
Therefore, the training labels do not uniquely entail the intended nonadjacent answer without a structural assumption. The independent checker confirms equal training predictions and different B,D predictions. This does not invalidate the task: spontaneous adoption of the intended structure is a legitimate behavioral capacity to test. It means that failed transfer can reflect failure to infer or deploy the hidden ordering, rather than absence of an ordinal computation once suitable premises are represented.\par
The same reasoning limits the assertion on p.7 that using a common mechanism across contexts would necessarily give similar performance. Encoding, memory, instruction comprehension, and elicitation can affect the observed outcome even when a reasoning component is available. The paper’s P4 and P5 explicitly recognize related behavior/mechanism and performance/competence distinctions. Its discussion also favors context-dependent capability statements. A claim about the observed protocol is supported more directly than a universal mechanism claim.\par
\subsection*{A symbolic-distance pattern does not uniquely identify TI}
For Experiment 1, consider an illustrative rule that memorizes the adjacent choices, always selects A if present in a novel pair, avoids G if present, and otherwise guesses. It has no general solution for interior novel comparisons. At a nonadjacent distance d, there are 7-d unordered pairs; the expected accuracy is (endpoint\_pairs + 0.5 * other\_pairs)/(7-d).\par
Distance | Total pairs | Endpoint-containing pairs | Expected accuracy\par
2 | 5 | 2 | 0.70\par
3 | 4 | 2 | 0.75\par
4 | 3 | 2 | 0.8333333333333334\par
5 | 2 | 2 | 1.00\par
6 | 1 | 1 | 1.00\par
The checker independently enumerates this curve. It demonstrates an alternative source of a positive distance association, without claiming that GPT-4o actually used this strategy or that the curve matches every plotted value. Endpoint-free and position-controlled comparisons would help discriminate mechanisms.\par
The authors describe behavior analogous to the symbolic distance effect; they do not prove a unique animal-like mechanism from the curve alone. All six fitted lines in Figure 2B point upward. Intermediate dips in Figure 2A do not contradict a positive fitted association. These are scope clarifications, not further confirmed errors.\par
\subsection*{Causality motivation}
The cited Halpern paper studies sufficient conditions for causal transitivity and explicitly discusses failures of transitivity in general. It does not establish ordinal TI as a necessary substrate of every causal judgment. The reviewed article’s motivation should be read with that qualification; it does not explicitly claim that every causal relation is transitive. This source check therefore does not refute the ordinal experiments or create another finding. [Halpern (2016), author PDF pp.2,6-8]\par
\section*{7. Versions, known issues, and overall dependence}
An immutable PMLR copy was independently downloaded by a support agent and matched the assigned PDF byte for byte. The author publication page links an older-title, 12-page draft bundle with an extra cover page and preliminary/under-review markers. Another author-hosted file is an 11-page anonymous preliminary draft. Neither was established as a later corrected edition. Exact upload dates were not established, and the drafts were not exhaustively diffed against the final article.\par
No relevant erratum or verified later repair was identified in the inspected primary records and targeted searches. OpenReview revision/discussion history remains inaccessible. The finite inconsistency is independently reproduced here, but no claim of priority or exhaustive absence of prior criticism is made.\par
The demonstrated problem affects reproducibility of one empirical control. The statistical and mechanistic limits restrict the interpretation of that case study. None demonstrates failure of the paper’s broader recommendation to use adversarial controls, stimulus variation, failure analysis, behavior/mechanism distinctions, and appropriate task accommodations.\par
\section*{8. Conclusion table}
Item | Judgment | Repair verification | Adoption | Remaining work\par
Endpoint exclusion versus distances 2-6 | Minor | Conditional 2-4 bound proved; actual implementation unresolved | Not requested | Obtain pair list and reconcile protocol/results\par
Chance, SD, intervals, and GLM | Needs Evidence | Published statistics not reproduced | Not requested | Raw responses, units, metadata, analysis specification\par
TI mechanism from transfer or SDE | Scope clarification | Countermodels checked; no claim they describe the real model | Not applicable | Discriminating controls and mechanistic evidence\par
Five-principle position | No Issue Found within inspected scope | Normative argument not experimentally certified | Not applicable | Human scientific assessment\par
\section*{9. Consolidated references}
[Rane-PMLR (2025)] Rane et al., PMLR 267:82051-82061. https://proceedings.mlr.press/v267/rane25a.html . Immutable matching PDF: https://raw.githubusercontent.com/mlresearch/v267/a3e4112da70e3f0623daeb3eeea06e2b34ed5e2f/assets/rane25a/rane25a.pdf .\par
[Halpern (2016)] Joseph Y. Halpern. Sufficient Conditions for Causality to Be Transitive. Philosophy of Science 83(2):213-226. DOI:10.1086/684915. Author PDF: https://www.cs.cornell.edu/home/halpern/papers/transitivity.pdf .\par
[Royka-publications] Author publication listing: https://amandaroyka.github.io/publications.html . Linked earlier bundle: https://amandaroyka.github.io/ICMLPOSITION.pdf . Additional preliminary draft inspected by support agent: https://amandaroyka.github.io/icml2025.pdf .\par
[OpenReview] Assigned record: https://openreview.net/forum?id=gCPJFcHskT . Earlier associated record: https://openreview.net/forum?id=RSbSX0Bf24 . These were discovered but their discussion and attachments were not accessible.\par
\section*{Supporting checks}
A support agent also inspected both figures independently.\par
Core expressions were:\par
{\ttfamily\footnotesize pairs = [(i,j,j-i) for i,j in combinations(range(1,6),2)\\
         if j-i > 1]\\
weights = [comb(n,j) for j in range(n+1)]\\
p\_two\_sided = sum(w for w in weights if w <= weights[k]) / 2**n}\par
Two initial script runs exited 1 because the output directory was absent. Between them, a replacement lacking explicit rewrite mode was rejected, and the unchanged script was inadvertently relaunched. File and process state were inspected before explicitly rewriting the script with directory creation. The corrected run succeeded.\par
The existing text extraction and image rendering provided coverage. No author program or model inference was executed; no packages were installed.\par
\section*{Supporting checks}
The original and supporting-download manifests preserve URLs, byte sizes and hashes. The saved earlier bundle has SHA-256 da90016eec08b6660920571167a1d6ee0daf38d3cb82c4467e7dfc4a528534d2. The Halpern PDF has SHA-256 5eed36f055bd20963c47e7a2f1d0f1c47009af8e474befd85724d3dd49b71306. These are supporting sources; the reviewed version remains the PMLR original.\par
Human scientific decisions and publication status remain pending.\par
\end{document}
